
\section{Evaluation Setup}

\cparagraph{Platforms.}
All experiments are conducted on a GPU workstation equipped with an NVIDIA RTX~3090 GPU (24\,GB memory, 936\,GB/s memory bandwidth) and a dual-socket Intel Xeon Gold 5218R CPU (80 logical cores).
We use PyTorch~2.5.1, Triton~3.1.0, CUDA~12.4, GCC~11.4.0, and Claude Sonnet~4~\cite{anthropic2025claude} as the underlying LLM.

\cparagraph{Benchmarks.}
We evaluate on three benchmark suites:
(1)~\emph{PolyBench/C}~\cite{polybench}: 30 numerical kernels (stencils, linear algebra, data mining) at default sizes ($N$=60--120) and 8$\times$ scale ($N$=480--960);
(2)~\emph{TSVC}~\cite{callahan1988tsvc}: 151 loop kernels covering linear traversals, reductions, indirect addressing, and loop-carried dependences;
(3)~8 application kernels from Rodinia~\cite{che2009rodinia} and ECP proxy applications~\cite{heroux2009mantevo}.

\cparagraph{Configurations.}
On PolyBench/C (our primary benchmark), we compare three configurations:
\begin{itemize}
\item \emph{Agent}: An autonomous LLM agent with tool use (write code, run tests, run benchmarks). The agent sees full error output and decides its own retry strategy. No compiler analysis. Up to 10 correctness turns + 3 optimization turns.
\item \emph{NA} (No Analysis): The LLM receives only the C source code with classified error feedback on retry. Same turn budget as the agent.
\item \emph{WA+Prof} (With Analysis + Profiling): Our full method. The LLM receives compiler-derived analysis facts. After correctness, NCU profiling feedback guides up to 3 optimization iterations.
\end{itemize}
On TSVC and application kernels, we report WA+Prof only (our method).

\cparagraph{Evaluation methodology.}
Speedups are measured over 50 iterations after 5 warmup iterations. Each iteration includes array re-initialization (GPU \texttt{clone()} for Triton, CPU \texttt{numpy.copy()} for the C reference) to ensure in-place kernels start from identical data. Triton timing includes \texttt{cuda.synchronize()}.
